% Glossary (AN-23 F-LOT: the term Heros Bold lower case on its own line,
% the definition below; bold words inside a definition are entries
% themselves).  CP-03: every \term of the book has an entry here and every
% entry is a \term of the book (make test-pubs-i123, i123_gate glossary).
% GATES: test-pubs-i123
\Gates{test-pubs-i123}
\BackChapter{Glossary}
This glossary defines the terms printed in bold where they are first used.

\GlossTerm{@ sign}{The character that begins an @function name; a cell entry
that begins with it is a formula.}
\GlossTerm{@function}{A built-in formula, such as \texttt{@SUM}, used as part
of a formula.  Chapter~13 describes them all.}
\GlossTerm{absolute reference}{A reference with a dollar sign before its
column letter and row number, \texttt{\$A\$1}, which names the same cell
wherever the formula is copied.}
\GlossTerm{add-in program}{A separate program loaded alongside Initech 123
that adds commands to it.}
\GlossTerm{address}{The column letter and row number that name a cell, such
as \texttt{B4}.}
\GlossTerm{argument separator}{The comma between the arguments of an
@function.}
\GlossTerm{arguments}{The values, cells or ranges an @function works on,
written in parentheses after its name.}
\GlossTerm{attach}{To load an add-in program into memory.}
\GlossTerm{borders}{Columns or rows printed on every page of a large print
range.}
\GlossTerm{cell}{The intersection of a column and a row of the worksheet;
it holds one entry.}
\GlossTerm{cell pointer}{The highlighted bar that marks the current cell.}
\GlossTerm{circular reference}{A formula that depends on its own value,
directly or through other formulas.  The \texttt{CIRC} indicator shows that
the worksheet contains one.}
\GlossTerm{control panel}{The top three lines of the Initech 123 window: the
current cell and its contents, the mode indicator, the entry being typed or
the menu, and the description of the highlighted menu item.}
\GlossTerm{criteria range}{A range that holds the conditions a record must
meet to be selected by a database @function.}
\GlossTerm{current cell}{The cell the cell pointer is in.}
\GlossTerm{current file}{The worksheet file last retrieved or saved;
\texttt{/File Save} offers its name.}
\GlossTerm{data range}{The range a Data or Graph command works on.}
\GlossTerm{database}{A table in the worksheet in which each row is a
record and each column a field.}
\GlossTerm{date serial number}{The number of days since 31 December 1899
by which Initech 123 stores a date: 1 January 1900 is 1.}
\GlossTerm{decimal places}{The number of digits a format shows after the
decimal point.}
\GlossTerm{field}{A column of a database.}
\GlossTerm{file name}{The name of a file on disk: up to eight characters,
optionally a period and an extension of up to three.}
\GlossTerm{footer}{A line of text printed at the bottom of every page.}
\GlossTerm{format}{The way a cell displays a number.  A format changes only
the display, never the number stored.}
\GlossTerm{formula}{An entry that calculates a value from numbers, cell
references and @functions.}
\GlossTerm{FROM range}{The range a copy or move is taken from.}
\GlossTerm{function name}{The name of an @function after the @ sign.}
\GlossTerm{global format}{The format of every cell not given one of its own.}
\GlossTerm{global setting}{A setting that applies to every cell without a
setting of its own: the global format, column width and label prefix.}
\GlossTerm{header}{A line of text printed at the top of every page.}
\GlossTerm{invoke}{To start an attached add-in program.}
\GlossTerm{key}{The field by which the records of a database are put in
order.}
\GlossTerm{label}{An entry of text.}
\GlossTerm{label prefix}{The character stored at the start of a label that
sets its alignment: \texttt{'} left, \texttt{"} right,
\texttt{\textasciicircum} centered, \texttt{\textbackslash} repeated.}
\GlossTerm{main menu}{The menu displayed when you press the slash key in
READY mode.}
\GlossTerm{margin}{The blank space between the edge of the paper and the
printed text.}
\GlossTerm{matrix}{A rectangular range of numbers treated as a unit.}
\GlossTerm{menu pointer}{The highlight on the menu item that ENTER selects.}
\GlossTerm{mixed reference}{A reference with a dollar sign before only its
column letter or only its row number, \texttt{\$A1} or \texttt{A\$1}.}
\GlossTerm{mode indicator}{The word at the right of the first line of the
control panel that shows what Initech 123 is doing: READY, LABEL, VALUE,
EDIT, MENU, FILES, RANGE or ERROR.}
\GlossTerm{named graph}{A set of graph settings kept under a name.}
\GlossTerm{natural order}{The recalculation order in which every formula is
calculated after the formulas it depends on.}
\GlossTerm{operating system}{The program that manages the files and runs the
programs of the computer: InitechDOS.}
\GlossTerm{page break}{A mark where one printed page ends and the next
begins.}
\GlossTerm{pie chart}{A graph that shows one data range as slices of a
circle.}
\GlossTerm{print range}{The range that is printed.}
\GlossTerm{prompt}{A question on the second line of the control panel.}
\GlossTerm{protected cell}{A cell that cannot be changed while worksheet
protection is enabled.  Cells are protected until they are unprotected.}
\GlossTerm{range}{A rectangle of cells, named by two opposite corners, such
as \texttt{A1..C5}, or by a range name.}
\GlossTerm{range name}{A name of up to 15 characters given to a range.}
\GlossTerm{recalculation order}{The order in which the formulas of the
worksheet are calculated: natural, columnwise or rowwise.}
\GlossTerm{record}{A row of a database.}
\GlossTerm{relative reference}{A reference without dollar signs,
\texttt{A1}, which names a cell by its distance from the formula.}
\GlossTerm{setup string}{Codes sent to a printer to select a mode or a
typeface.}
\GlossTerm{status line}{The bottom line of the Initech 123 window: messages
and the status indicators.}
\GlossTerm{string}{Text used in a formula: text in double quotation marks, a
label in a cell, or the result of a string @function.  The \texttt{\&}
operator joins two strings.}
\GlossTerm{TO range}{The place a copy or move goes, named by its upper left
cell.}
\GlossTerm{unsaved work}{Changes to the worksheet since it was last saved or
retrieved, and an entry being typed.}
\GlossTerm{value}{A number or a formula.}
\GlossTerm{worksheet}{The grid of 256 columns and 8,192 rows in which you
work.}
\GlossTerm{worksheet area}{The part of the window that shows the cells.}
\GlossTerm{worksheet file}{A file on disk that holds a worksheet, in the
\texttt{.WK1} format.}
\GlossTerm{worksheet frame}{The column letters and row numbers around the
worksheet area.}
\GlossTerm{X range}{The range that supplies the labels along the horizontal
axis of a graph.}
